\section{Framework Design}\label{design}

\subsection{Architecture and Lifecycle Layers}
The framework is built around three lifecycle layers, each with its own scope. Figure~\ref{fig:framework-layers} shows where they meet the existing services.

\begin{figure}[tbp]
  \centering
    \scalebox{0.95}{\begin{tikzpicture}[
    font=\small,
    box/.style={
        draw, rounded corners, align=center,
        minimum width=7.0cm, minimum height=0.95cm, inner sep=3pt
      },
    iacbox/.style={
        draw=user, thick, rounded corners, align=center,
        minimum width=7.0cm, minimum height=0.95cm, inner sep=3pt,
        fill=user!10
      },
    drift/.style={-Latex, thick},
    note/.style={align=left, font=\footnotesize},
    framebox/.style={
        draw=user, very thick, rounded corners=8pt,
        inner xsep=10pt, inner ysep=10pt
      },
    resolvenote/.style={align=left, font=\footnotesize, text width=2.8cm},
    resolvearrow/.style={-Latex, thick, draw=user},
    node distance=0.55cm
  ]

  \node[box] (svc) {Service / Application layer\\(\footnotesize team-specific procedures, overrides)};
  \node[box, below=of svc] (cm) {Configuration Management layer\\(\footnotesize Puppet + Foreman)};
  \node[box, below=of cm] (cloud) {Virtual Infrastructure layer\\(\footnotesize OpenStack resources, metadata, lifecycle)};
  \node[box, below=of cloud] (phys) {Physical Infrastructure layer\\(\footnotesize hardware inventory, network registration, lifecycle records)};

  \node[framebox, fit=(svc)(cm)(cloud)(phys)] (framework) {};

  \node[iacbox, above=0.6cm of framework.north] (iac) {\textbf{IaC integration layer}\\(\footnotesize Terraform/OpenTofu + CI/CD)};

  \draw[drift, draw=user] (iac.south) -- (framework.north);

  \coordinate (b1) at ($(svc.south)!0.5!(cm.north)$);
  \coordinate (b2) at ($(cm.south)!0.5!(cloud.north)$);
  \coordinate (b3) at ($(cloud.south)!0.5!(phys.north)$);

  \draw[drift] (svc.south) -- (cm.north);
  \draw[drift] (cm.south) -- (cloud.north);
  \draw[drift] (cloud.south) -- (phys.north);

  \node[resolvenote, right=0.3cm of framework.east |- b1] (d1)
  {declarative handover:\\bootstrap + classification\\derived from code};
  \draw[resolvearrow] (d1.west) -- (framework.east |- b1);

  \node[resolvenote, right=0.3cm of framework.east |- b2] (d2)
  {aligned lifecycle:\\provisioning + classification\\from single definition};
  \draw[resolvearrow] (d2.west) -- (framework.east |- b2);

  \node[resolvenote, right=0.3cm of framework.east |- b3] (d3)
  {unified registration:\\custom providers for\\inventory + certificates};
  \draw[resolvearrow] (d3.west) -- (framework.east |- b3);

  \node[note, left=0.3cm of framework.west |- cm] (sot)
  {single\\source of\\intent (VCS)};
  \draw[drift, draw=user] (sot.east) -- (framework.west |- cm);

\end{tikzpicture}
}
  \caption{View of the proposed IaC framework. Encapsulation of configuration, provisioning and registration through an integration layer.}
  \label{fig:framework-layers}
\end{figure}

The virtual infrastructure layer (OpenStack resources, metadata, lifecycle) and the physical infrastructure layer (hardware inventory, network registration, lifecycle records) are both managed through Terraform/OpenTofu. Custom providers take care of resource provisioning, inventory registration, and certificate management in one place, which removes the manual and error-prone steps that these two layers used to require.

The CM layer keeps full responsibility for runtime configuration through Puppet and Foreman. The framework does not replace or duplicate them; it treats them as the source of truth for anything at host level.

What matters in this design is that intent is declared once and stored in VCS, and that provisioning, classification, and registration are all derived from it. That is what addresses the scattered state described in Section~\ref{infra}, and it removes drift at each of the integration points in Figure~\ref{fig:framework-layers}.

\subsection{State and Drift Model}

Terraform/OpenTofu detects drift by comparing what was declared against what actually exists. Running a plan reconciles three things: the declared configuration, the stored state file, and the real resource state queried through the provider APIs. Anything that disagrees is surfaced before a change is applied.

We use this by running plan operations on a schedule through CI/CD. Drift introduced outside the framework, typically a manual intervention during an incident, shows up at the next pipeline run, so operators get an explicit and auditable list of what no longer matches. Fixing it means re-applying the declared state through the normal workflow.

\subsection{Provisioning Workflow}

Provisioning runs through CI/CD. An operator commits a change to their infrastructure repository. The pipeline validates it and runs a plan, which gives an auditable view of what is about to change. Once approved, the pipeline applies it. Terraform/OpenTofu provisions the cloud resources and registers the host with the supporting services through the custom providers. The instance boots, cloud-init runs the bootstrap logic, and the Puppet agent makes its first catalog request. From there, configuration convergence belongs to CM and the framework is out of the picture. Figure~\ref{fig:provisioning-seq} shows the sequence.

\begin{figure}[tbp]
  \centering
  \scalebox{0.95}{\begin{tikzpicture}[
  font=\footnotesize,
  block/.style={draw, rounded corners, align=center, minimum height=0.8cm, minimum width=1.4cm},
  arr/.style={-Latex, thick},
  lbl/.style={above=12pt, font=\scriptsize, align=center}
]

\node[block] (op) {Operator};
\node[block, right=0.9cm of op] (git) {Git};
\node[block, right=0.9cm of git] (ci) {CI/CD};
\node[block, right=0.9cm of ci] (tf) {Terraform};
\node[block, right=0.7cm of tf] (infra) {OpenStack +\\CERN services};
\node[block, right=0.7cm of infra] (puppet) {Puppet};

\draw[arr] (op) -- node[lbl]{commit} (git);
\draw[arr] (git) -- node[lbl]{trigger} (ci);
\draw[arr] (ci) -- node[lbl]{plan + apply} (tf);
\draw[arr] (tf) -- node[lbl]{provision +\\register} (infra);
\draw[arr, dashed] (infra) -- node[lbl]{bootstrap +\\handover} (puppet);

\end{tikzpicture}
}
  \caption{Provisioning workflow: infrastructure intent flows from Git through CI/CD and Terraform to infrastructure services. After bootstrap, Puppet assumes runtime configuration.}
  \label{fig:provisioning-seq}
\end{figure}

The point of the design is that provisioning and CM inputs come from one declarative definition, which is what removes the manual coordination at the provisioning/CM boundary. Updates go through the same \texttt{plan/apply} cycle.
